\documentclass[10pt]{article}
\usepackage[letterpaper,margin=0.9in]{geometry}
\usepackage{amsmath,amsthm}
\usepackage{newtxtext,newtxmath}
\theoremstyle{definition}
\newtheorem{theorem}{Theorem}
\usepackage{booktabs,tabularx,array}
\usepackage{float}
\usepackage{microtype}
\usepackage[hidelinks]{hyperref}
\setlength{\parindent}{0pt}
\setlength{\parskip}{2.5pt}
\setlength{\abovedisplayskip}{5pt}
\setlength{\belowdisplayskip}{5pt}
\setlength{\textfloatsep}{7pt}
\setlength{\floatsep}{6pt}
\setlength{\intextsep}{6pt}
\renewcommand{\arraystretch}{1.08}
\newcommand{\Vbar}{\bar V}
\newcommand{\Dbar}{\bar D}
\newcommand{\shat}{\widehat s}
\newcommand{\R}{\mathbb R}
\newcommand{\T}{\mathsf T}

\title{\vspace{-1.4em}\bfseries Confidence Transport for Adaptive Experimentation:\\Theory and a Controlled Audit\vspace{-0.35em}}
\author{Bahram Behzadian\textsuperscript{1}\quad
Khurshid Juarev\textsuperscript{2}\quad
Diego Granziol\textsuperscript{3}\quad
Houssam Nassif\textsuperscript{1}\\[-0.2em]
{\small \textsuperscript{1}Meta \quad
\textsuperscript{2}Purestrength AI \quad
\textsuperscript{3}University of Oxford}}
\date{}

\begin{document}
\maketitle
\vspace{-2.1em}

\begin{abstract}
Model updates change the sensitivities used for allocation, while
reward confidence may rely on sensitivities fixed when observations
were collected. We give a finite-action confidence-transfer result
under explicit pre-reward metric and solver certificates, yielding
reward optimism and a regret bound on their joint validity event.
Under fixed-reference and approximation-envelope assumptions, a
corrected center supplies confidence without nonlinear optimizer convergence.
In a nearly linear 29-parameter audit, all primary-policy float64 confidence
and optimism diagnostics pass, yet the evaluated bound is vacuous and the
multiplier policy has higher sample-mean pseudo-regret than its omission
within each primary-horizon condition.
\end{abstract}
\vspace{-0.4em}

\section{Adaptive model updates and confidence transport}
Adaptive allocation may retrain a reward model, then select an ad or ranking action from predicted rewards and uncertainty; related neural-bandit methods use learned representations \cite{zhou2020neural,riquelme2018deep}. Retraining changes replayed sensitivities and the Gram matrix. Collection-time confidence does not directly bound scores in the recomputed metric.

At round $t\le T$, observe $x_t$, choose $a_t$ from a fixed finite
nonempty set $\mathcal A$, then observe $r_t$; the confidence target is
the conditional mean reward $\mu_t^*(a)$ used for allocation, not a
treatment-effect estimand. Queries $q_t(a)\in\R^d$ have fixed coordinates
and dimension. Let $\mathcal H_t^-$ contain the context and history before
round-$t$ randomization and $\mathcal H_t$ contain all pre-reward
information. The metrics $\Vbar_t,V_t$, radius $\beta_t\ge0$, and
$\theta_t$ when used are $\mathcal H_t^-$-measurable (predictable);
other score inputs and the selected action are finite and
$\mathcal H_t$-measurable. All validity events are measurable.

For $\lambda,\sigma>0$, with $I$ the $d$-dimensional identity, define reference geometry and widths by
\begin{equation}
\Vbar_1=\lambda I,\qquad
\Vbar_{t+1}=\Vbar_t+\sigma^{-2}q_t(a_t)q_t(a_t)^\top,
\qquad
\bar s_t(a)^2=q_t(a)^\top\Vbar_t^{-1}q_t(a).
\label{eq:reference}
\end{equation}
Define cumulative pseudo-regret and frozen information gain, with $\gamma_0=0$, as
\begin{equation}
R_T=\sum_{t=1}^T\left[\max_{a\in\mathcal A}\mu_t^*(a)-\mu_t^*(a_t)\right],\qquad
\gamma_T=\log\frac{\det\Vbar_{T+1}}{\det\Vbar_1}.
\label{eq:regret-gain}
\end{equation}
Let $V_t$ be the exact current metric and $C_t$ the algorithmic operator, both $d\times d$ symmetric positive-definite (SPD). Set $s_t(a)^2=q_t(a)^\top V_t^{-1}q_t(a)$ and $s_{C,t}(a)^2=q_t(a)^\top C_t^{-1}q_t(a)$.

\begin{theorem}[Finite-action confidence transfer]\label{thm:transfer}
Under this protocol, suppose $\|q_t(a)\|_2\le G<\infty$.
Define $E_{\rm conf}$ as the event that, simultaneously for
$t\le T$ and $a\in\mathcal A$,
\begin{equation}
|\mu_t^*(a)-m_t(a)|
\le\beta_t\bar s_t(a)+b_t(a),\qquad b_t(a)\ge0.
\label{eq:refconf}
\end{equation}
Let $\Dbar_t\ge0$,
$0<\kappa_{-,t}\le\kappa_{+,t}<\infty$, and
$1\le\alpha_t<\infty$ be finite pre-reward factors.
Define $E_{\rm cert}$ as simultaneous validity, for every $t\le T$,
of the Loewner comparisons
\begin{equation}
e^{-\Dbar_t}\Vbar_t\preceq V_t\preceq e^{\Dbar_t}\Vbar_t,\qquad
\kappa_{-,t}V_t\preceq C_t\preceq\kappa_{+,t}V_t,
\label{eq:certificates}
\end{equation}
and, for a nonnegative pre-reward width map,
\begin{equation}
s_{C,t}(a)\le\shat_t(a)\quad(a\in\mathcal A),\qquad
\shat_t(a_t)\le\alpha_t s_{C,t}(a_t).
\label{eq:solver}
\end{equation}
Use this same realized width map in the score and selector:
\begin{equation}
U_t(a)=m_t(a)+\beta_t e^{\Dbar_t/2}
\sqrt{\kappa_{+,t}}\,\shat_t(a)+b_t(a),\qquad
U_t(a_t)\ge\max_{a\in\mathcal A}U_t(a)-\xi_t,\quad\xi_t\ge0.
\label{eq:score}
\end{equation}
On $E_{\rm conf}\cap E_{\rm cert}$,
$U_t(a)\ge\mu_t^*(a)$ for every action and round, and pathwise
\begin{equation}
\begin{aligned}
R_T\le{}&
\sqrt{\left(\sigma^2+\frac{G^2}{\lambda}\right)\gamma_T
\sum_{t=1}^T\beta_t^2
\left[1+\alpha_t e^{\Dbar_t}
\sqrt{\frac{\kappa_{+,t}}{\kappa_{-,t}}}\right]^2}\\
&+2\sum_{t=1}^T b_t(a_t)+\sum_{t=1}^T\xi_t.
\end{aligned}
\label{eq:regret}
\end{equation}
If $\Pr(E_{\rm conf})\ge1-\delta$ and
$\Pr(E_{\rm cert})\ge1-\delta_{\rm cert}$, these conclusions hold
with probability at least $1-\delta-\delta_{\rm cert}$.
\end{theorem}
\emph{Proof sketch.} On the joint event, reference confidence, the inverse
forms of \eqref{eq:certificates}, and all-action solver upper validity make
the score optimistic. Approximate maximization, lower confidence at the played
action, the reverse comparisons, and played-action sharpness give the
roundwise coefficient in \eqref{eq:regret}. Summation, Cauchy--Schwarz,
and the frozen rank-one potential bound
$\sum_t\bar s_t(a_t)^2\le(\sigma^2+G^2/\lambda)\gamma_T$
yield the result. Thompson geometry supplies endpoint and path-length
constructions for the first sandwich \cite{nussbaum1994finsler}; Section~2
supplies \eqref{eq:refconf} without optimizer convergence.

\section{A corrected center without optimizer convergence}
Use a differentiable model mean $\mu_\theta(x,a)$ in fixed parameter coordinates $\theta\in\Theta\subseteq\R^d$, with $\Theta$ a convex smoothness region. For $z=(x,a)$, define $q_\theta(z)=\nabla_\theta\mu_\theta(z)$, $q_t(a)=q_{\theta_t}(x_t,a)$, $z_s=(x_s,a_s)$, and $q_s:=q_{\theta_s}(x_s,a_s)$. The representation path $\theta_t\in\Theta$ is predictable. Fix before data collection $\theta^\circ\in\Theta$, $\|\theta^\circ\|_2\le S$, with $\mu_t^*(a)=\mu^*(x_t,a)$ and $\mu^*(z)=\mu_{\theta^\circ}(z)+m^\circ(z)$. Observed rewards obey $r_t=\mu_t^*(a_t)+\eta_t$ and $\mathbb E[e^{u\eta_t}\mid\mathcal H_t]\le e^{u^2\sigma^2/2}$ for every real $u$.

Supply nonnegative pre-reward bounds $\epsilon_{\rm lin}(s)$ on the absolute Taylor remainder of $\mu_{\theta^\circ}(z_s)$ about $\theta_s$, and $\epsilon_{\rm mis}(s)$ for $|m^\circ(z_s)|$. Define the current actionwise envelopes $\epsilon_{{\rm lin},t}(a),\epsilon_{{\rm mis},t}(a)$ analogously about $\theta_t$; the past sum below is known before selection. Set
\begin{align}
y_s&=r_s-\mu_{\theta_s}(z_s)+q_s^\top\theta_s, &
\widehat\theta_t^{\rm lin}&=\Vbar_t^{-1}\sigma^{-2}\sum_{s<t}q_sy_s, \\
m_t^{\rm corr}(a)&=\mu_{\theta_t}(x_t,a)+q_t(a)^\top(\widehat\theta_t^{\rm lin}-\theta_t).&&
\label{eq:center}
\end{align}
For $\delta\in(0,1)$, self-normalization \cite{abbasi2011improved} and the Taylor envelopes give \eqref{eq:refconf} with probability at least $1-\delta$, simultaneously over rounds and actions, using $m_t=m_t^{\rm corr}$,
\begin{equation}
\beta_t^{\rm corr}=\sqrt{\gamma_{t-1}+2\log(1/\delta)}+\sqrt{\lambda}S
+\frac1\sigma\sqrt{\sum_{s<t}[\epsilon_{\rm lin}(s)+\epsilon_{\rm mis}(s)]^2},
\quad b_t(a)=\epsilon_{{\rm lin},t}(a)+\epsilon_{{\rm mis},t}(a).
\label{eq:beta}
\end{equation}
The exact current squared-loss generalized Gauss-Newton (GGN) metric is $V_t=\lambda I+\sigma^{-2}\sum_{s<t}q_{\theta_t}(z_s)q_{\theta_t}(z_s)^\top$.

\section{Controlled audit}
We ran four policies in a nearly linear scaled-tanh bandit with five actions and 29 parameters. For normalized features $\|\phi(z)\|_2\le1$, the mean is $\mu_{\theta,W}(z)=\sqrt W\tanh(\phi(z)^\top\theta/\sqrt W)$. Exact realizability holds with $\|\theta^\circ\|_2\le R=1$; parameters are projected to radius $R$, $\lambda=1$, and Gaussian noise has $\sigma=0.25$. For each horizon $T\in\{250,500,1000\}$ and target $D_{\rm target}\in\{0.25,0.5,1,2\}$, the predeclared rule sets
\begin{equation}
W=\left[\frac{4c_hR\sqrt{T-1}}{\sigma\sqrt\lambda D_{\rm target}}\right]^2,
\qquad c_h=\frac4{3\sqrt3},\qquad L_g=\frac{c_h}{\sqrt W},\qquad
Q_t=\sum_{s<t}\|\theta_t-\theta_s\|_2^2,\quad D_{Q,t}=\frac{2L_g\sqrt{Q_t}}{\sigma\sqrt\lambda}.
\label{eq:experiment-scale}
\end{equation}
For $u=\phi(z)^\top\theta$, $|u|\le1$, and
$f_W(u)=\sqrt W\tanh(u/\sqrt W)$,
$0\le1-f_W(u)/u\le u^2/(3W)\le1/(3W)$, using the continuous extension
at zero; at $T=1000$, the bound $1/(3W)$ ranges from
$1.375\times10^{-7}$ to $8.798\times10^{-6}$ across the four conditions.
Here $L_g$ bounds the operator norm of the mean Hessian along replay segments. $D_{\rm target}$ equals $D_{Q,T}$ when every replay displacement is maximal ($2R$). Four policies, three separately run horizons, four targets, and 50 evaluation seeds give 2,400 trajectories. Ten disjoint tuning seeds selected shared optimizer settings by mean squared prediction error. Policies share contexts and per-action noise draws; their histories diverge. All use dense float64 Cholesky solves and exact argmax over the five actions, with failure level $\delta=0.05$. The exact-arithmetic primary instantiation has $C_t=V_t$, $\kappa_{-,t}=\kappa_{+,t}=\alpha_t=1$, and $\xi_t=0$.

The transport-endpoint policy is a dense endpoint-certificate comparator. It uses
$d_{{\rm Th},t}=\max_i|\log\lambda_i(V_t,\Vbar_t)|$, with
$\lambda_i$ the generalized eigenvalues of $(V_t,\Vbar_t)$:
the smallest common exponent for the first symmetric sandwich in
\eqref{eq:certificates} \cite{nussbaum1994finsler}.
Both matrices are known before the reward, so the certificate is valid
in exact arithmetic; its float64 value is not a verified enclosure.
Each policy scores
$m_t^{\rm corr}(a)+\beta_t^{\rm corr}w_t(a)+b_t(a)$ on its own history:
Widths are $w_t=e^{D_{Q,t}/2}s_t$ (transport Hessian),
$e^{d_{{\rm Th},t}/2}s_t$ (transport endpoint),
$\bar s_t$ (frozen reference), and $s_t$ (naive current).
The primary policy sets $\Dbar_t=D_{Q,t}$; ``Hessian'' names the
envelope argument. Let $\mathcal D_{{\rm selected},t}$ be the exact
Thompson-Finsler length of the SPD path that interpolates each replay
parameter from $\theta_s$ to $\theta_t$; then
$d_{{\rm Th},t}\le\mathcal D_{{\rm selected},t}\le D_{Q,t}$, and
$\mathcal D_{{\rm quad},t}$ is its 32-point Gauss-Legendre estimate
at the declared checkpoints.

\begin{table}[H]
\centering\footnotesize
\begin{tabular}{@{}lrrrr@{}}
\toprule
$D_{\rm target}$ & $0.25$ & $0.5$ & $1$ & $2$ \\
\midrule
Transport Hessian mean $R_{1000}$ & 81.05 & 82.44 & 87.04 & 93.99 \\
Transport endpoint mean $R_{1000}$ & 79.05 & 79.49 & 80.06 & 81.89 \\
Frozen reference mean $R_{1000}$ & 79.09 & 80.13 & 80.25 & 81.67 \\
Naive current mean $R_{1000}$ & 79.00 & 79.49 & 80.10 & 81.88 \\
\midrule
Median $D_Q/\mathcal D_{\rm quad}$ & $1.88\times10^6$ & $9.77\times10^5$ & $4.82\times10^5$ & $2.51\times10^5$ \\
Median trajectory-maximum $d_{\rm Th}$ & $1.31\times10^{-7}$ & $5.23\times10^{-7}$ & $2.10\times10^{-6}$ & $8.74\times10^{-6}$ \\
Median $e^{D_Q/2}$ & 1.02 & 1.04 & 1.09 & 1.18 \\
\bottomrule
\end{tabular}
\caption{Primary horizon $T=1000$. Regret means average 50 trajectories. The last three rows concern the primary policy: path ratios pool eligible seed-checkpoint pairs with $\mathcal D_{\rm quad}>10^{-12}$; endpoint distances are maximized within each trajectory before taking the median; inflation medians pool seeds and all rounds, not terminal rounds or run maxima.}
\label{tab:audit}
\end{table}

Within each $T=1000$ condition, the analytic multiplier raised sample
mean pseudo-regret relative to the otherwise matched naive-current
policy (Table~\ref{tab:audit}); histories diverge after the first
differing action, and the contrast is descriptive:
$D_{\rm target}$ also changes $W$, so it is not a dose response. In each of the 12 horizon-target cells, 50/50 primary trajectories passed both all-action, all-round reference-confidence and optimism diagnostics, with cellwise exact 95\% Clopper-Pearson interval $[92.9\%,100\%]$. Base seeds are reused across cells. These tolerance-based float64 checks and quadrature values are diagnostics, not verified enclosures. Across the 12 primary cells, medians of per-trajectory ratios of the numerical right-hand side (RHS) of \eqref{eq:regret} to pseudo-regret range from $126.7$ to $155.7$. Every primary trajectory's terminal RHS exceeds the $2T$ pseudo-regret cap from $|\mu^*(z)|\le1$.

\section{Implications for adaptive experimentation}
On the joint confidence-and-certificate event, the theorem gives optimism and
a regret bound, but valid certificates need not improve allocation: in this
nearly linear simulation, the analytic multiplier accompanies higher sample-mean
pseudo-regret than its omission within each primary-horizon condition.
The audit does not establish a uniform policy ordering or treatment-effect
inference after adaptive assignment.

\vspace{-0.25em}
\begin{thebibliography}{9}\setlength{\itemsep}{0pt}\setlength{\parskip}{0pt}\small
\bibitem{abbasi2011improved} Y. Abbasi-Yadkori, D. P\'al, and C. Szepesv\'ari. Improved Algorithms for Linear Stochastic Bandits. \emph{NeurIPS}, 2011.
\bibitem{zhou2020neural} D. Zhou, L. Li, and Q. Gu. Neural Contextual Bandits with UCB-based Exploration. \emph{ICML}, 2020.
\bibitem{riquelme2018deep} C. Riquelme, G. Tucker, and J. Snoek. Deep Bayesian Bandits Showdown: An Empirical Comparison of Bayesian Deep Networks for Thompson Sampling. \emph{ICLR}, 2018.
\bibitem{nussbaum1994finsler} R. D. Nussbaum. Finsler structures for the part metric and Hilbert's projective metric and applications to ordinary differential equations. \emph{Differential and Integral Equations}, 7:1649--1707, 1994.
\end{thebibliography}
\end{document}
